% from ega4-2-fr.tex
\oldpage[IV-2]{78}

\begin{proposition}[4.7.6]
\label{IV.2.4.7.6-fr}
Soit
\[
0\longrightarrow\mathscr F'\longrightarrow\mathscr F
\longrightarrow\mathscr F''\longrightarrow0
\]
une suite exacte de $\mathcal O_X$-Modules cohérents. Si en un point $x\in X$
la longueur géométrique $\lambda_x(\mathscr F)$ de $\mathscr F$ en ce point
est définie, il en est de même des longueurs géométriques
$\lambda_x(\mathscr F')$ et $\lambda_x(\mathscr F'')$, de $\mathscr F'$ et
$\mathscr F''$ au point $x$, et réciproquement; en outre, on a
\begin{equation}
\lambda_x(\mathscr F)=\lambda_x(\mathscr F')+\lambda_x(\mathscr F'').
\tag{4.7.6.1}
\label{IV.2.4.7.6.1-fr}
\end{equation}
\end{proposition}

\begin{proof}
La proposition résulte aussitôt de la définition, car
$\operatorname{long}(\mathscr F_x)$ est finie si et seulement si
$\operatorname{long}(\mathscr F'_x)$ et
$\operatorname{long}(\mathscr F''_x)$ sont finies, et l'on a
\[
\operatorname{long}(\mathscr F_x)
=\operatorname{long}(\mathscr F'_x)+\operatorname{long}(\mathscr F''_x).
\]
\end{proof}

\begin{env}[4.7.7]
\label{IV.2.4.7.7-fr}
Sous les hypothèses de (4.7.5), les seuls points $x\in X$ tels que
$\mathscr F_x$ soit un $\mathcal O_x$-module de longueur finie non nulle sont
(en vertu de (3.1.2)) les points maximaux de
$\operatorname{Supp}(\mathscr F)$, comme il résulte de Bourbaki,
\emph{Alg. comm.}, chap.~IV, \S~2, n\textsuperscript{o}~5, cor.~2 de la
prop.~7, la question étant locale sur $X$. On sait qu'il existe un
sous-préschéma fermé $Y$ ayant $\operatorname{Supp}(\mathscr F)$ pour espace
sous-jacent, et un $\mathcal O_Y$-Module cohérent $\mathscr G$ tel que
$\mathscr F=j_*(\mathscr G)$, $j:Y\to X$ étant l'injection canonique. Si $x$
est point maximal de $Y$, on voit que les longueurs géométriques de
$\mathscr F$ et de $\mathscr G$ en $x$ sont égales, $\mathcal O_{X,x}$ et
$\mathcal O_{Y,x}$ ayant même corps résiduels. On en déduit la caractérisation
suivante de la longueur géométrique:
\end{env}

\begin{proposition}[4.7.8]
\label{IV.2.4.7.8-fr}
Sous les hypothèses de (4.7.5), soit $x$ un point maximal de
$\operatorname{Supp}(\mathscr F)$. Soit $k'$ une extension parfaite de $k$;
posons $X'=X\otimes_k k'$, $\mathscr F'=\mathscr F\otimes_k k'$; alors la
longueur géométrique $\lambda_x(\mathscr F)$ de $\mathscr F$ en $x$ est égale
à la longueur du $\mathcal O_{x'}$-module $\mathscr F'_{x'}$, en tout point
maximal $x'$ de $\operatorname{Supp}(\mathscr F')$ au-dessus de $x$. En outre,
il existe une extension radicielle finie $k_1$ de $k$ telle que, en posant
$X_1=X\otimes_k k_1$, $\mathscr F_1=\mathscr F\otimes_k k_1$,
$\lambda_x(\mathscr F)$ soit égale à la longueur du
$\mathcal O_{x_1}$-module $(\mathscr F_1)_{x_1}$ en tout point maximal $x_1$
de $\operatorname{Supp}(\mathscr F_1)$ au-dessus de $x$.
\end{proposition}

\begin{proof}
La remarque de (4.7.7) montre qu'on peut se borner au cas où $x$ est point
maximal de $X$. En outre, la question est locale sur $X$ et l'on peut donc
supposer $X=\operatorname{Spec}(A)$ affine noethérien,
$\mathscr F=\widetilde M$, où $M$ est un $A$-module de type fini; soient
$B=A\otimes_k k'$, $\mathfrak p$ (resp. $\mathfrak q$) l'idéal premier
minimal de $A$ (resp. $B$) correspondant à $x$ (resp. $x'$). On a alors, en
appliquant la formule (4.7.2.1)
\[
\operatorname{long}_{B_{\mathfrak q}}(\mathscr F'_{x'})
=\operatorname{long}_{A_{\mathfrak p}}(\mathscr F_x)
 \operatorname{long}_{B_{\mathfrak q}}
 (B_{\mathfrak q}/\mathfrak pB_{\mathfrak q}).
\]
Mais comme $k(x)=A_{\mathfrak p}/\mathfrak pA_{\mathfrak p}$, le même
raisonnement que dans (4.7.3) montre que
$\operatorname{long}_{B_{\mathfrak q}}
(B_{\mathfrak q}/\mathfrak pB_{\mathfrak q})$ est égal à la longueur de
l'anneau local $(B/\mathfrak pB)_{\mathfrak q'}$, où
$\mathfrak q'=\mathfrak q/\mathfrak pB$. L'anneau
$(B/\mathfrak pB)_{\mathfrak q'}$ est un anneau local de
$(A/\mathfrak p)\otimes_k k'$ pour un idéal premier minimal de cet anneau,
donc aussi un anneau local de l'anneau total des fractions de
$k(x)\otimes_k k'$, et la longueur d'un tel anneau est par définition
$\mu_i(k(x)\mid k)$, d'où la première assertion. La seconde se démontre de
même, en remplaçant $k'$ par $k_1$ et utilisant (4.7.3, (iii)).
\end{proof}

On dit encore, dans le cas où $x$ est point maximal de
$\operatorname{Supp}(\mathscr F)$, que la longueur géométrique de $\mathscr F$
en $x$ est la multiplicité radicielle du cycle premier maximal
$\overline{\{x\}}$ associé à $\mathscr F$.

\begin{corollary}[4.7.9]
\label{IV.2.4.7.9-fr}
Sous les hypothèses de (4.7.5), soient $k'$ une extension de $k$,
$X'=X\otimes_k k'$, $\mathscr F'=\mathscr F\otimes_k k'$; soient $Z$ une
composante irréductible de $\operatorname{Supp}(\mathscr F)$, et $Z'$ une
composante irréductible de $\operatorname{Supp}(\mathscr F')$ qui domine $Z$
(4.2.7). Alors les multiplicités radicielles de $Z$ par rapport à
$\mathscr F$ et de $Z'$ par rapport à $\mathscr F'$ sont les mêmes.
\end{corollary}

